\begin{tikzpicture}[scale=1]

% === Parâmetros ===
\def\s{2}   % raio do círculo ao quadrado
\def\a{1}   % coordenada dos interseptos

\pgfmathsetmacro{\r}{sqrt(\s)} % raio do círculo

  \begin{axis}[      
      xlabel={$x$}, ylabel={$y$},
      axis equal image,        % mantém escala x=y para círculos não ficarem ovalados
      xmin=-1.2*\r, xmax=1.2*\r, ymin=-1.2*\r, ymax=1.2*\r,
      xtick={-\a,0,\a,\r}, ytick={-\a,0,\a},
      xticklabels={-\a,0,\a,$\sqrt{\s}$}, 
      x tick label style={anchor=south west}, % rótulos abaixo e a esquerda do eixo X
      y tick label style={anchor=south east}, % rótulos acima e a esquerda do eixo y
      % width=8cm, height=6cm,
      every axis x label/.style={at={(ticklabel* cs:0.97)}, anchor=north},
      every axis y label/.style={at={(ticklabel* cs:0.97)}, anchor=east},
    ]

    % --- Linhas Tracejadas ---
     \addplot[Tracejado] coordinates {(-\a,0) (-\a,\a)};
     \addplot[Tracejado] coordinates {(-\a,\a) (0,\a)};
     \addplot[Tracejado] coordinates {(\a,0) (\a,-\a)};
     \addplot[Tracejado] coordinates {(0,-\a) (\a,-\a)};

    % --- Curvas ---

    \draw[thick, blue] (axis cs:0,0) circle [radius=\r];    % Círculo de raio r

    \addplot[name path=reta, domain=-1.1*\r:1.1*\r, samples=100, thick, red] {-x}node[pos=1, below, red] {$y=-x$};

    % --- Região sombreada com Arc ---
    \path[fill=gray!30, opacity=0.6]
      (axis cs:\a,-\a)     % Ponto inicial, 
      % necessário para o comando Arc a seguir.
      % Se vc adicionar um "--" depois do ponto acima, 
      % o "arc" não vai saber de onde começar
      arc[start angle=-45, end angle=135, radius=\r] % Arco do ciclo
      -- cycle;          % Conecta o ponto inicial ao final

    % --- Pontos de interseção ---
    \addplot[only marks, black!60, mark=*] coordinates {(-\a,\a) (\a,-\a) };


  \end{axis}
  
\end{tikzpicture}
